\documentclass[11pt,a4paper]{article}
\usepackage[utf8]{inputenc}
\usepackage[T1]{fontenc}
\usepackage[spanish]{babel}
\usepackage{geometry}
\usepackage{hyperref}
\usepackage{enumitem}
\usepackage{tabularx}
\usepackage{xcolor}

\geometry{margin=2cm}
\hypersetup{colorlinks=true, linkcolor=blue, urlcolor=blue}

\title{Cambios de Infraestructura — Sesión 2026-05-09}
\author{Rimay Guide — Frontend}
\date{}

\begin{document}
\maketitle

\section{Stack instalado}
\begin{itemize}[nosep]
  \item \texttt{zustand} v5.0.13 — state management (1~KB, sin Provider)
  \item \texttt{vite-plugin-pwa} v1.3.0 — service worker + manifest automáticos
  \item \texttt{react-router} v7.13.0 — routing declarativo (ya instalado, ahora usado)
  \item \texttt{@supabase/supabase-js} v2.105.4 — cliente inicializado (sin backend aún)
  \item \texttt{leaflet} v1.9.4 + \texttt{react-leaflet} v5.0.0 — mapa real OpenStreetMap
\end{itemize}

\section{Estructura de carpetas}
\begin{verbatim}
src/
  stores/          authStore, tourStore, locationStore (Zustand)
  hooks/           useGeolocation (watchPosition + Haversine)
  workers/         downloadWorker (Cache API con postMessage)
  services/        tourService (fetchTourBySlug mock)
  lib/             supabaseClient (cliente inicializado)
  features/        auth/, tour/ (preparado para escalar)
\end{verbatim}

\section{Arquitectura de decisiones}

\begin{tabularx}{\textwidth}{l|X|X}
  \textbf{Decisión} & \textbf{Elegido} & \textbf{Justificación} \\
  \hline
  State mgmt   & Zustand     & Stores por dominio, sin Provider, 1~KB. Escala con middlewares (persist, devtools). \\
  Routing      & react-router & URLs con slug necesarias para QR codes. Soporta lazy loading. \\
  PWA          & vite-plugin-pwa & Genera SW + manifest en build. Estrategias: SWR (imágenes), CacheFirst (MP3), NetworkFirst (API). \\
  Mapa         & Leaflet + OSM & Gratis, sin API key, probado desde 2011. \\
  Geo hook     & Custom (refs) & Fórmula Haversine + distance-to-stop. Refs rompen loop infinito de dependencias. \\
  Worker       & Web Worker nativo & Sin deps. Migrable a Comlink si escala. \\
\end{tabularx}

\section{Flujo de usuario}
\begin{center}
  \texttt{/login} $\rightarrow$ \texttt{/} (SplashScreen con stops inline) $\rightarrow$ \texttt{/player?stopId=X}
\end{center}
Todas las rutas protegidas (\texttt{/}, \texttt{/player}, \texttt{/tour/:slug}) redirigen a \texttt{/login} si no hay sesión activa. El parámetro \texttt{?redirect=} preserva la ruta de destino.

\section{Stores (Zustand)}
\begin{itemize}[nosep]
  \item \textbf{authStore}: \texttt{user}, \texttt{isAuthenticated}, \texttt{login(email, pass)}, \texttt{socialLogin(provider)}, \texttt{logout()}.
  \item \textbf{tourStore}: \texttt{tour}, \texttt{stops}, \texttt{currentStopIndex}, \texttt{isDownloaded}, \texttt{downloadProgress}.
  \item \textbf{locationStore}: \texttt{position} (lat, lon, accuracy, timestamp), \texttt{activeStopId}, \texttt{isWatching}.
\end{itemize}

\section{useGeolocation hook}
\begin{itemize}[nosep]
  \item Usa \texttt{navigator.geolocation.watchPosition} con alta precisión.
  \item Distancia Haversine entre posición y cada parada (umbral configurable, default 5~s).
  \item \textbf{Bug corregido}: loop infinito por \texttt{stops} en array de dependencias del \texttt{useEffect}. Solución: \texttt{useRef} actualizados vía suscripción de Zustand que retorna \texttt{null}.
\end{itemize}

\section{LocationModal (Leaflet)}
\begin{itemize}[nosep]
  \item Mapa con tiles de OpenStreetMap.
  \item 9 marcadores con coordenadas reales de Sacsayhuamán ($-13.507^\circ$, $-71.981^\circ$).
  \item Color por estado: terracota (actual), verde (completado), gris (futuro).
  \item Posición del usuario en tiempo real con círculo de precisión.
  \item GPS se activa manualmente (botón ``Activar GPS'' por privacidad).
  \item Centrado automático del mapa al obtener la primera posición.
\end{itemize}

\section{DownloadWorker}
\begin{itemize}[nosep]
  \item Recibe array de URLs + cacheName vía \texttt{postMessage}.
  \item Descarga cada archivo con \texttt{fetch} y lo guarda en Cache API.
  \item Reporta progreso (\texttt{current}, \texttt{total}, \texttt{percent}) al hilo principal.
  \item Soporta cancelación vía mensaje \texttt{cancel}.
\end{itemize}

\section{PWA — Service Worker}
\begin{itemize}[nosep]
  \item Generado por \texttt{vite-plugin-pwa} (workbox-based).
  \item \textbf{App Shell}: precache de JS, CSS, HTML.
  \item \textbf{Imágenes} (Unsplash): Stale-While-Revalidate (30 días, 50 entradas).
  \item \textbf{Audios MP3}: CacheFirst (90 días, 100 entradas).
  \item \textbf{API}: NetworkFirst con timeout de 5~s (1 hora, 100 entradas).
  \item \texttt{registerType: autoUpdate} — actualiza el SW sin pedir confirmación.
\end{itemize}

\section{Qué hace la aplicación actualmente}
\begin{itemize}[nosep]
  \item Login con email/contraseña y Google/Apple (UI funcional, lógica mock).
  \item SplashScreen con hero image, lista de 9 paradas inline, botón ``Ver ubicación''.
  \item Navegación completa entre pantallas vía react-router.
  \item Protección de rutas: si no hay sesión, redirige a \texttt{/login}.
  \item AudioPlayer con reproducción de MP3, waveform SVG, seek, skip $\pm$15s, next/prev.
  \item Mapa Leaflet con OpenStreetMap mostrando 9 paradas georreferenciadas en Sacsayhuamán.
  \item Geolocalización en tiempo real con marcador de usuario y círculo de precisión.
  \item Descarga offline (DownloadModal con UI + Web Worker para Cache API).
  \item PWA funcional: service worker, manifest, instalable en pantalla de inicio.
  \item Soporte para URL directa vía QR: \texttt{/tour/sacsayhuaman}.
\end{itemize}

\section{Qué no hace}
\begin{itemize}[nosep]
  \item No hay backend real. Los datos de tours son mock (hardcodeados en App.tsx).
  \item No hay autenticación real. El login simula éxito sin validar contra Supabase ni ningún servidor.
  \item No hay descarga real de archivos. El Web Worker está implementado pero el DownloadModal usa una barra de progreso simulada (setInterval).
  \item No hay reproducción automática al entrar en rango de una parada. El hook useGeolocation detecta proximidad pero no dispara el cambio de stop en el player.
  \item No hay persistencia de estado entre sesiones (zustand sin middleware persist).
  \item No hay dark mode. Las variables están definidas en theme.css pero no hay toggle ni lógica.
  \item No hay tests de ningún tipo.
  \item No hay animaciones de transición entre pantallas.
  \item No hay manejo de errores de red (offline detection, retry).
  \item Solo 1 audio real (sacsayhuaman\_es.mp3). Los otros 8 son placeholders.
\end{itemize}

\section{Qué falta por hacer (próximos pasos)}
\begin{enumerate}[nosep]
  \item \textbf{Conectar Supabase}: reemplazar mocks en \texttt{authStore} y \texttt{tourService} por llamadas reales a Supabase Auth y PostgreSQL.
  \item \textbf{Descarga real}: cablear el \texttt{downloadWorker} con el \texttt{DownloadModal} para descargar archivos MP3 de Supabase Storage a Cache API con progreso real.
  \item \textbf{Trigger de audio por geolocalización}: cuando \texttt{useGeolocation} detecta que el usuario entró en el radio de una parada, navegar automáticamente a \texttt{/player?stopId=X} y reproducir.
  \item \textbf{Manejo de offline}: mostrar banner ``Sin conexión'' cuando no hay red, usar datos cacheados (SW) para la UI, degradar gracefulmente.
  \item \textbf{Persistencia de sesión}: middleware \texttt{persist} de Zustand para guardar auth token, tour descargado y stops completados en localStorage/IndexedDB.
  \item \textbf{Dark mode}: implementar toggle con \texttt{next-themes} (ya instalado) y adaptar componentes a variables \texttt{.dark}.
  \item \textbf{Transiciones}: animar cambios de ruta con \texttt{motion} (framer-motion, ya instalado) — slide, fade entre pantallas.
  \item \textbf{Audios restantes}: producir y subir los 8 archivos MP3 faltantes para completar el tour de Sacsayhuamán.
  \item \textbf{Tests}: configurar Vitest + React Testing Library, testear stores, hooks y componentes críticos (auth, player, geolocalización).
  \item \textbf{CI/CD}: GitHub Actions para lint, typecheck, test y deploy a Vercel/Netlify.
  \item \textbf{Accesibilidad}: revisar contraste, focus traps en modales, \texttt{aria-*} labels, soporte lector de pantalla.
  \item \textbf{Optimización}: lazy loading de rutas (\texttt{React.lazy}), code splitting por feature, análisis de bundle.
\end{enumerate}

\section{Correcciones}
\begin{itemize}[nosep]
  \item \textbf{tsconfig.json}: creado con \texttt{strict: true}, paths \texttt{@/*}, soporte WebWorker.
  \item \textbf{Colores}: LoginScreen migrado de \texttt{\#1a365d} a \texttt{var(--terracotta)}.
  \item \textbf{AudioPlayer}: \texttt{isPlaying} arranca en \texttt{false} (sin auto-reproducción). Botón ``Volver'' agregado.
  \item \textbf{LoginRoute/SplashRoute}: ahora ambas exigen autenticación.
  \item \textbf{.env}: variables \texttt{VITE\_SUPABASE\_URL} y \texttt{VITE\_SUPABASE\_ANON\_KEY} placeholder.
\end{itemize}

\vfill
\hrule
\small \textbf{Archivos nuevos}: 13 (stores, hooks, workers, services, lib, tsconfig, .env). \\
\textbf{Archivos modificados}: 7 (App.tsx, main.tsx, vite.config.ts, SplashScreen, LoginScreen, AudioPlayer, TourStopsList, index.css). \\
\textbf{Build}: exitoso, 0 errores. Bundle final: 461~KB JS + 120~KB CSS (gzip: 142~KB + 23~KB).

\end{document}
